\cotefile{91}
\scanpage{99}
% L1
\[
= \varinjlim_{\unsure{\mathfrak P|\mathcal P}}\ \varprojlim_j\
\deleted{\unsure{\operatorname{Hom}}}\prod_\alpha
\operatorname{Hom}(X_\alpha^{\mathfrak P},Y_j)
\]

\begin{center}\rule{1.4cm}{0.4pt}\end{center}

% L2--5; vertical pen stroke beside the heading and the opening of the sentence.
\noindent\llap{\raisebox{-1.1ex}{\rule{0.35pt}{2.4ex}}\hspace{.35em}}Ensembles quasi-schématiques.

Les ens quasi-schématiques de type fini sur un \deleted{\unsure{schéma}} ens quasi-schématique \uline{noethérien} forment une \uline{catégorie géométrique} (?)
% archive sheet number: 98
